% !TEX program = lualatex
\documentclass[11t]{article}

\usepackage{microtype}
\usepackage{geometry}
%\geometry{margin=1in}
\usepackage{setspace}
\usepackage{hyperref}
\usepackage[style=apa, backend=biber, natbib=true]{biblatex}
\addbibresource{references.bib}

\usepackage{amsmath, amssymb}
\usepackage{booktabs}
\usepackage{longtable}

\usepackage{fontspec}
\setmainfont{Carlito}

\usepackage[scaled=.95]{helvet}

\hypersetup{
  colorlinks=true,
  linkcolor=blue,
  citecolor=blue,
  urlcolor=blue
}

\setlength{\parindent}{0pt}
\setlength{\parskip}{0.7em}
\onehalfspacing

\newcommand{\MVS}{\mathrm{MVS}^{\ast}}
\newcommand{\sigmoid}{\operatorname{sigmoid}}
\newcommand{\logonep}{\operatorname{log1p}}
\newcommand{\Indic}{\mathbb{I}}

\title{\textbf{TECHNICAL NOTE: The MVS Framework}\\
Underwriting External Credit via Digitized ROSCA Behaviour}
\author{}
\date{16 February 2026}
\begin{document}
\maketitle

\textbf{Subject:} Operationalizing ROSCA Behaviour for Initial Credit Scoring and PD Calibration

\vspace{0.5em}
\hrule
\vspace{0.8em}

\section{The narrative: the information problem and why ROSCAs matter}

The central constraint to extending credit in financially thin environments is information asymmetry. When bureau histories are incomplete or absent, lenders cannot directly observe repayment type. 
Rotating Savings and Credit Associations (ROSCAs)—tontines, osusus—matter because they partially solve this information problem through repeated interaction, peer observability, and embedded discipline.
The theoretical foundation is clear: ROSCAs can be understood as repeated-game institutions with endogenous enforcement \citet{BesleyCoateLoury1993}. Empirically, ROSCAs often function as commitment devices that help members maintain disciplined periodic payments \citet{Gugerty2007}. 
Governance quality, sanctions, and social connectedness shape enforcement and sustainability over time \citet{McNabbLeMayBoucherBonan2019}. In bidding ROSCAs, the bid discount behaves like an implicit interest rate and reveals liquidity stress and shocks \citet{Klonner2003,CalomirisRajaraman1998}.
A key bridge to formal underwriting is the insight that \emph{publicly observable characteristics} inside these informal institutions can predict creditworthiness. 
\citet{SchreinerNagarajan2000} provide direct evidence from Gambian ASCRAs/ROSCAs that cheap, observable reliability markers contain predictive signal for credit screening. 
This is precisely the operational premise of MVS: turn socially embedded observables into an interpretable underwriting signal at launch, then calibrate it into a formal probability of default once outcomes accumulate.

\section{Literature-grounded mechanisms: what we are measuring}

The score is intentionally mechanism-based rather than purely correlational. We map observed ROSCA behaviour to economic channels identified in the literature. Payment discipline and persistence capture time-consistent commitment \citet{Gugerty2007}. Allocation order and post-payout slippage encode incentive compatibility and moral hazard after receiving the pot \citet{BesleyCoateLoury1993}. Governance and enforcement strength reflect heterogeneity in rules, sanctions, and connectedness \citet{McNabbLeMayBoucherBonan2019}. In bidding ROSCAs, discount level and volatility reveal liquidity stress in the upper tail \citet{Klonner2003,CalomirisRajaraman1998}. Finally, social capital proxies (repeat interaction, centrality, endorsements) raise reputational stakes and strengthen informal enforcement \citet{SchreinerNagarajan2000}.

\section{Target, label, and decision-time discipline (anti-leakage)}

The long-run objective is to estimate PD@3M: the probability that a loan reaches DPD60 within 90 days after disbursement. We define the calibration label:
\begin{equation}
\text{Def3M}_i = \Indic\{\text{DPD60 occurs on or before day 90 post-disbursement}\}.
\end{equation}

All predictors must be computed using only information available at the underwriting decision timestamp $T_i$. This ``decision-time discipline'' prevents leakage and ensures deployability.

Because payout timing matters, define:
\begin{equation}
\text{AllocByDec}_i = \Indic\{\text{AllocDate}_i \le T_i\}.
\end{equation}
Post-payout features are valid only when $\text{AllocByDec}_i=1$.

\section{Data dictionary (minimum viable schema)}

To keep the system auditable and partner-ready, we define a minimum viable schema. Variable names are intentionally short for operational simplicity. All timestamps are required to respect the decision-time cutoff.

\begin{longtable}{p{0.18\textwidth}p{0.12\textwidth}p{0.12\textwidth}p{0.50\textwidth}}
\toprule
\textbf{Field} & \textbf{Level} & \textbf{Type} & \textbf{Definition}\\
\midrule
\endfirsthead
\toprule
\textbf{Field} & \textbf{Level} & \textbf{Type} & \textbf{Definition}\\
\midrule
\endhead
mid & member & string & Member identifier \\
gid & group & string & ROSCA/group identifier \\
cid & cycle & string & Cycle identifier \\
mdate & meeting & date & Monthly meeting date \\
t\_dec & decision & datetime & Underwriting decision timestamp (feature cutoff) \\
due & meeting & number & Amount due \\
paid & meeting & number & Amount paid \\
dlate & meeting & int & Days late (0 if on time) \\
ont & meeting & \{0,1\} & On-time flag: $1$ if dlate$=0$ \\
N & cycle & int & Group size \\
aord & member/cycle & int & Allocation order position (1..N) \\
adate & member & date & Allocation date (pot receipt) \\
rtype & group & enum & ROSCA type: RANDOM / BIDDING / COMMITTEE \\
rules & group & \{0,1\} & Documented rules exist and communicated \\
san6 & member/group & int & Count of sanctions in last 6 meetings \\
santype & meeting & enum & LATE\_FEE / ABSENCE\_FINE / WARNING / SUSPENSION / OTHER \\
sanamt & meeting & number & Monetary amount of sanction \\
sure\_n & member & int & Number of sureties (0/1/2+) \\
sure\_type & member & enum & NONE / INDIVIDUAL / POOLED \\
sure\_str & member & enum & NONE / WEAK / STRONG \\
bid\_flag & attempt & \{0,1\} & Bid attempt made \\
disc & attempt & number & Bidding discount (\%) recorded per attempt \\
rep & member & \{0,1\} & Repeat participation in same group across cycles \\
cent & member & \{0,1\} & Centrality above group median \\
end\_f & member & \{0,1\} & Foreman endorsement (clean) \\
end\_s & member & \{0,1\} & Senior endorsement (clean) \\
loanid & loan & string & Loan identifier \\
t\_disb & loan & datetime & Disbursement timestamp \\
ten & loan & int & Tenor (days) \\
freq & loan & enum & WEEKLY / BIWEEKLY / MONTHLY \\
dpd60 & loan & \{0,1\} & DPD60 flag \\
def3m & loan & \{0,1\} & $\text{Def3M}$ label (defined above) \\
\bottomrule
\end{longtable}

\section{From mechanisms to measurement (narrative \texorpdfstring{$\rightarrow$}{->} guided math)}

The design objective of MVS is to remain interpretable while reflecting the non-linear nature of behavioural risk. Early history is disproportionately informative relative to incremental history; 
tail volatility and clustered slippage matter more than small average differences. We therefore use smooth, monotone non-linear maps—sigmoid, exponential decay, and log transforms—so that the score saturates where it should and becomes sharply sensitive where risk concentrates.

\subsection{Eligible window and time decay}

Let $K_i$ be the number of eligible monthly meetings prior to $T_i$ (standard track: $K_i \ge 6$). Index meetings $m=1,\dots,K_i$ (oldest $\rightarrow$ newest). Define time-decay weights:
\begin{equation}
w_{im}=\lambda^{K_i-m}, \qquad 0<\lambda<1,
\end{equation}
with a conservative initial setting $\lambda=0.80$ for monthly cadence (sensitivity-tested during calibration).

\subsection{Time-decayed behavioural primitives}

From meeting-level primitives, compute time-decayed on-time rate and average lateness:
\begin{equation}
W_i = \sum_{m=1}^{K_i} w_{im}.
\end{equation}

\begin{equation}
\text{OTR}_i = \frac{1}{W_i}\sum_{m=1}^{K_i} w_{im}\,\text{ont}_{im},
\qquad
\text{AL}_i  = \frac{1}{W_i}\sum_{m=1}^{K_i} w_{im}\,\max(0,\text{dlate}_{im}).
\end{equation}
Let $\text{LS}_i$ be the maximum consecutive late streak in the eligible window.

We use:
\begin{equation}
\sigmoid(x)=\frac{1}{1+e^{-x}}, \qquad \logonep(x)=\ln(1+x).
\end{equation}

\subsection{Pillar 1: Payment discipline (PDIS)}

Map primitives into smooth monotone components:
\begin{align}
S^{OTR}_i &= 20\cdot \sigmoid\!\left(k_{OTR}(\text{OTR}_i-c_{OTR})\right),\\
S^{AL}_i  &= 6\cdot \exp\!\left(-a_{AL}\cdot \logonep(\text{AL}_i)\right),\\
S^{LS}_i  &= 9\cdot \exp\!\left(-a_{LS}\cdot \text{LS}_i\right).
\end{align}
A small recovery term $S^{RC}_i$ (``rapid cure'') can be included when late/missed payments are cured within 7 days.

Normalize to a 35-point scale:
\begin{equation}
S^{PDIS}_i = \frac{35}{37}\left(S^{OTR}_i+S^{AL}_i+S^{LS}_i+S^{RC}_i\right).
\end{equation}

\subsection{Pillar 2: Order and post-payout incentives (ORDR)}

Define allocation bucket $\text{AB}_i \in \{\text{EARLY},\text{MID},\text{LATE}\}$ using terciles of $\text{aord}/N$. Assign base multipliers:
\begin{equation}
b(\text{LATE})=1.0,\quad b(\text{MID})=0.6,\quad b(\text{EARLY})=0.3.
\end{equation}
Define post-payout slippage $\text{Slip}_i$ as an indicator equal to 1 if a new late streak of length $\ge 2$ begins after allocation and before $T_i$ (valid only when $\text{AllocByDec}_i=1$). Then:
\begin{equation}
S^{ORDR}_i = 15\cdot b(\text{AB}_i)\cdot (1-\pi_{Slip}\cdot \text{Slip}_i).
\end{equation}

\subsection{Pillar 3: Governance and enforcement (GOV)}

Map rules and sanctions via monotone functions:
\begin{align}
S^{Rules}_i &= 5\cdot \sigmoid\!\left(k_R(\text{rules}_i-0.5)\right),\\
S^{San}_i   &= 6\cdot \exp(-a_Z\cdot \text{san6}_i).
\end{align}
Let $S^{Sure}_i \in \{0,3,6\}$ map surety strength NONE/WEAK/STRONG. Normalize:
\begin{equation}
S^{GOV}_i = \frac{20}{17}\left(S^{Rules}_i+S^{San}_i+S^{Sure}_i\right).
\end{equation}

\subsection{Pillar 4: Liquidity stress in bidding ROSCAs (LIQ)}

For bidding ROSCAs, let $A_i$ be bid attempts in the eligible window. Define attempt volatility via interquartile range:
\begin{equation}
\text{IQR}_i = Q_{75}(\text{disc}_{ia}) - Q_{25}(\text{disc}_{ia}), \quad a\in A_i.
\end{equation}
Let $q_i\in[0,1]$ be the member's discount quantile rank within the group attempt distribution. Define:
\begin{align}
S^{Lvl}_i &= 6\cdot \left(1-\sigmoid\!\left(k_q(q_i-q_0)\right)\right),\\
S^{Vol}_i &= 6\cdot \exp\!\left(-a_v\cdot \logonep(\text{IQR}_i/v_{ref})\right),\\
S^{LIQ}_i &= \frac{15}{12}\left(S^{Lvl}_i+S^{Vol}_i\right).
\end{align}

\subsection{Pillar 5: Social capital (SOC)}

Use a bounded linear combination of cheap observables:
\begin{equation}
\text{socraw}_i = 4\,\text{rep}_i + 3\,\text{cent}_i + 2\,\text{end\_f}_i + 2\,\text{end\_s}_i,
\qquad
S^{SOC}_i = \frac{15}{11}\,\text{socraw}_i.
\end{equation}

\subsection{Final score and guardrails}

Compute the raw total:
\begin{equation}
\MVS^{raw}_i = S^{PDIS}_i + S^{ORDR}_i + S^{GOV}_i + S^{LIQ}_i + S^{SOC}_i.
\end{equation}

Guardrails are applied last to prevent tail-risk dilution and structural blind spots. Decision bands (pre-calibration): Green ($\MVS \ge 75$), Amber ($60 \le \MVS \le 74$), Red ($\MVS < 60$).

\section{Strengths, limits, and operational use}

MVS is interpretable: each pillar corresponds to a literature-backed mechanism, enabling transparent reason codes for approvals/declines. It also addresses cold start by enabling underwriting without formal bureau history using repeated behavioural data.

At launch, MVS is not a calibrated PD; it is a decisioning and ranking tool. Accept-only bias is a key risk if lending occurs only to high-score applicants. These limitations are addressed through a governed Phase 1 pilot.

\section{The calibration imperative (phases and why Phase 1 must be governed)}

Stable PD calibration requires outcome data and sufficient default events; event counts, not only loan counts, determine model stability \citet{RileyEtAl2020,VittinghoffMcCulloch2006}. Overly complex models too early risk overfitting. The appropriate approach is phased:

Phase 0 (data readiness): lock taxonomy, schema, timestamps, QA rules, and ensure MVS can be computed deterministically.

Phase 1 (go-live with controlled exploration): deploy MVS with guardrails and reason codes; approve conservatively while allocating a bounded learning budget to generate outcomes beyond the acceptance frontier.

Phase 2 (first calibration): fit parsimonious logistic and/or survival models mapping MVS (and a small set of primitives) to Def3M; validate out-of-time and by group clusters (gid).

Phase 3 (scale and sophistication): as event volume grows, move to monotone GAMs and a monotone GBM challenger; refresh calibration regularly and monitor drift. Probability calibration should be monitored explicitly (see \citealp{SklearnCalib} for standard calibration diagnostics).

\section{Phase 1 pilot governance: numeric limits a partner can sign}

To bound downside while enabling learnable outcomes, the first 500 loans should run under explicit numeric governance constraints. The Core sleeve is at least 85\% of approvals (Green and high-Amber). The Exploration sleeve is capped at 15\% (low-Amber and a small, selected subset of high-potential Red profiles).

Let $E_0$ be the standard initial exposure cap agreed with the partner. In Phase 1, per-borrower maximum exposure is capped at $E_0$ in Core and $0.35E_0$ in Exploration. Products should generate delinquency signals within 90 days: tenor 90--120 days with weekly or biweekly installments.

Hard stop-loss triggers are applied mechanically. Pause new Exploration approvals if Exploration DPD30 exceeds 12\% over the most recent 60 disbursement days, or if Exploration DPD60 exceeds 6\%. Pause all new approvals if portfolio-level DPD60 exceeds 4\% over the most recent 60 disbursement days. Outcomes are monitored by weekly vintages, and each decision stores pillar-level reason codes and guardrail flags.

\section{Strategic positioning (partner discussion)}

This approach is not merely a formula; it is an engineered learning loop. We begin with a mechanism-based, auditable score grounded in ROSCA theory and evidence \citet{BesleyCoateLoury1993,Gugerty2007,McNabbLeMayBoucherBonan2019,Klonner2003,CalomirisRajaraman1998,SchreinerNagarajan2000}. We go live conservatively with guardrails and reason codes, run a governed pilot with bounded exploration, and then calibrate PD@3M with parsimonious models before scaling, consistent with event-driven methodology \citet{RileyEtAl2020,VittinghoffMcCulloch2006}.

\clearpage
\section*{References}
\nocite{*}
\printbibliography
\end{document}
